Logical anomaly detectors are judged by whether they flag a faulty scene and how well their maps rank faulty regions. Neither number tells an inspector how many intact parts a deployed map will send them to. We build matched rearrangements of 138 real MVTec LOCO pushpin trays: the same intact pin is moved within its compartment (a valid sham) or into an occupied neighbor (a fault), and a vacancy-only control removes it. SALAD's composition branch detects every fault and ranks the affected compartments at 98.2\% AUROC. Yet at a cutoff calibrated for 95\% target recall it opens 11.1 of 13 intact compartments per faulty tray, because inserting a pin raises 90\% of the untouched compartments, while removing one or moving it validly does not. We trace this spread to how SALAD labels synthetic insertions: over the whole inserted class. In controlled retraining with identical inputs, labelling only the pasted footprint removes the spread (to under 2\% of compartments). On the 91 native logical pushpin anomalies it more than doubles localization, raising AU-sPRO at 5\% false-positive rate from 34\% to 72--79\%, and from 16\% to 70--84\% on extra-pin defects. In screw bags, whose extra parts are interchangeable and annotated as such, it brings nothing. A label should point where the task locates the fault. The same matched controls show that two appearance detectors, which look better under fault-only localization, alarm on about 90\% of valid shams, and that SALAD's full system misses its nominal false-alarm rate on normal test trays.
